% ============================================================
% РАЗДЕЛ 3. Промежуточное решение на базе NFS
% ============================================================
\section{ПРОМЕЖУТОЧНОЕ РЕШЕНИЕ НА БАЗЕ NFS И КООРДИНАЦИИ ЗАГРУЗОК}

В данном разделе описывается промежуточный этап оптимизации системы бессерверного инференса: переход от изолированной загрузки весов модели в том \texttt{emptyDir} к использованию сетевой файловой системы NFS в качестве общего хранилища. Раскрываются недостатки исходной схемы, механизм интеграции NFS в Kubernetes через ресурсы PersistentVolume и PersistentVolumeClaim, проблемы синхронизации параллельных загрузок, особенности файловых блокировок поверх сетевых файловых систем, применение механизма Leader Election через Kubernetes Lease API, а также системные ограничения, обусловившие необходимость разработки распределённой P2P/FUSE-архитектуры.

\subsection{Исходная схема: storage-initializer и emptyDir}

\subsubsection*{Паттерн init-контейнера и жизненный цикл тома emptyDir}

Исходная схема развёртывания модели в Kubernetes использует паттерн init-контейнера. Под (pod) содержит два контейнера: \texttt{storage-initializer}, запускаемый первым в режиме init, и runtime-контейнер, являющийся основным. Init-контейнер монтирует том типа \texttt{emptyDir}, запускает процесс загрузки весов модели из внешнего репозитория (HuggingFace, ModelScope, объектное S3-совместимое хранилище) и записывает файлы в смонтированный том. После успешного завершения init-контейнера Kubernetes запускает runtime-контейнер (vLLM, TGI, Triton), который читает веса из того же тома \texttt{emptyDir} и загружает их в GPU-память для обслуживания запросов.

Том типа \texttt{emptyDir} создаётся Kubernetes в момент планирования пода на узел и физически размещается на локальном диске этого узла, либо, при явной конфигурации, в памяти. Ключевая характеристика данного типа тома состоит в том, что его жизненный цикл жёстко привязан к жизненному циклу пода: при удалении или перепланировании пода том уничтожается вместе со всем содержимым. Том \texttt{emptyDir} не является персистентным и принципиально не разделяется между подами: каждый под получает собственную изолированную директорию.

\subsubsection*{Фундаментальный недостаток: повторные загрузки}

Отсутствие разделяемого хранилища порождает фундаментальный недостаток данной схемы. При создании новой реплики, даже на той же узловой машине, каждый под выполняет полную загрузку файлов модели независимо от всех ранее запущенных подов. Если уже работающий под уже завершил загрузку той же модели, новый под этого не знает и не может воспользоваться имеющейся копией данных.

При масштабировании до $N$ реплик внешний репозиторий получает $N$ независимых последовательных или параллельных запросов на один и тот же набор файлов. Для современных моделей класса LLM с объёмом весов от 7 до 70~ГБ это означает:
\begin{itemize}[leftmargin=*]
  \item суммарный исходящий трафик на уровне десятков и сотен гигабайт при каждом событии масштабирования;
  \item время cold start каждой реплики полностью определяется скоростью загрузки из внешнего источника, а не скоростью чтения с диска;
  \item нагрузка на внешний канал связи возрастает пропорционально числу реплик, что может приводить к деградации пропускной способности и к превышению rate-лимитов репозитория.
\end{itemize}

Таким образом, схема \texttt{storage-initializer + emptyDir} не решает задачу cold start при масштабировании: повторные загрузки по-прежнему присутствуют, хотя и перемещены с уровня повторных запусков того же пода на уровень каждой новой реплики.

\subsection{Переход к NFS как общему хранилищу}

\subsubsection*{NFS и режим ReadWriteMany}

Network File System (NFS) --- зрелый протокол сетевого доступа к файлам, разработанный компанией Sun Microsystems и стандартизированный в версиях NFSv3 (RFC~1813) и NFSv4 (RFC~7530). В контексте Kubernetes NFS привлекателен прежде всего тем, что поддерживает режим доступа \texttt{ReadWriteMany} (RWX), при котором несколько подов на разных узлах кластера могут одновременно монтировать один и тот же том как для чтения, так и для записи \cite{nfs_k8s_docs}. Большинство блочных хранилищ поддерживают лишь режим \texttt{ReadWriteOnce} (RWO), допускающий монтирование только одним подом, что исключает их применение для общего кэша весов моделей.

\subsubsection*{Интеграция в Kubernetes: PV и PVC}

Схема использования NFS в Kubernetes предполагает создание двух ресурсов: \texttt{PersistentVolume} (PV), описывающего физическое хранилище, и \texttt{PersistentVolumeClaim} (PVC), представляющего запрос пода на использование этого хранилища. Ниже приведены примеры манифестов.

\begin{lstlisting}[language=yaml, caption={PersistentVolume на базе NFS}]
apiVersion: v1
kind: PersistentVolume
metadata:
  name: models-nfs-pv
spec:
  capacity:
    storage: 500Gi
  volumeMode: Filesystem
  accessModes:
    - ReadWriteMany
  persistentVolumeReclaimPolicy: Retain
  nfs:
    path: /exports/models
    server: 10.0.1.100
\end{lstlisting}

\begin{lstlisting}[language=yaml, caption={PersistentVolumeClaim для NFS-тома}]
apiVersion: v1
kind: PersistentVolumeClaim
metadata:
  name: models-nfs-pvc
  namespace: inference
spec:
  accessModes:
    - ReadWriteMany
  resources:
    requests:
      storage: 500Gi
  volumeName: models-nfs-pv
\end{lstlisting}

Поды, смонтировавшие этот PVC через \texttt{spec.volumes}, получают доступ к единой файловой системе, физически расположенной на NFS-сервере. Kubernetes автоматически связывает PVC с подходящим PV при наличии совпадающих параметров ёмкости и режима доступа. Для автоматизации создания PV по запросу PVC без ручного администрирования применяется динамическое выделение ресурсов (dynamic provisioning) через внешние провизионеры, такие как \textit{NFS Subdir External Provisioner}, которые автоматически создают поддиректории на NFS-сервере для каждого нового PVC.

\subsubsection*{Устранение повторных загрузок}

Введение NFS принципиально изменяет характер загрузок: первый под, обнаруживший отсутствие модели в NFS-директории, загружает её из внешнего репозитория и записывает в общее хранилище. Все последующие поды, монтирующие тот же PVC, обнаруживают файлы уже готовыми и пропускают этап загрузки. Таким образом, при масштабировании с 1 до $N$ реплик загрузка из внешнего источника выполняется ровно один раз, а не $N$ раз, что устраняет ключевую неэффективность схемы \texttt{emptyDir}.

\subsection{Проблема параллельного скачивания и условие гонки}

\subsubsection*{Race condition при первом развёртывании}

Несмотря на устранение повторных загрузок при масштабировании уже существующей модели, переход к NFS порождает новую задачу: состояние гонки (race condition) при первом развёртывании или после принудительной инвалидации кэша. В этом сценарии несколько реплик запускаются одновременно. Каждая из них на этапе инициализации проверяет директорию модели в NFS и обнаруживает её отсутствие. После этого несколько реплик независимо начинают параллельную загрузку весов по одному и тому же пути назначения.

Параллельная запись одних и тех же файлов из нескольких источников приводит к непредсказуемым последствиям: файлы могут оказаться повреждёнными вследствие чередования блоков от разных источников, директория модели может находиться в частично записанном состоянии, которое воспринимается другими подами как готовое. В результате runtime-контейнер может попытаться загрузить модель из повреждённого набора файлов, что приведёт к ошибке инициализации или некорректным ответам модели.

\subsubsection*{Алгоритм координированной загрузки}

Для решения задачи консистентности при параллельном старте разработан алгоритм координированной загрузки, основанный на сочетании файловых блокировок и атомарных файловых операций. Алгоритм реализует паттерн «один загрузчик --- множество ожидающих»:

\begin{enumerate}[leftmargin=*]
  \item Под проверяет наличие финального маркер-файла \texttt{.ready} в директории модели. Если маркер присутствует, модель готова к использованию; загрузка не требуется, выполнение переходит к запуску runtime-контейнера.
  \item Если маркер отсутствует, под открывает специальный файл блокировки \texttt{.lock} в директории модели и выполняет попытку получения эксклюзивной блокировки (\texttt{LOCK\_EX | LOCK\_NB}).
  \item Под, успешно получивший блокировку, становится лидером загрузки. Перед началом загрузки лидер повторно проверяет наличие маркера \texttt{.ready} — другой лидер мог успеть завершить загрузку в промежутке между проверкой и получением блокировки. При обнаружении маркера лидер освобождает блокировку и переходит к запуску runtime-контейнера.
  \item Лидер начинает загрузку весов модели во временную директорию \texttt{<model>\_tmp}, физически расположенную в том же NFS-разделе. Запись ведётся во временную директорию, что гарантирует отсутствие частично записанных данных по целевому пути.
  \item После завершения загрузки всех файлов лидер выполняет верификацию целостности: сравнивает суммарный объём файлов и, при наличии манифеста, проверяет контрольные суммы.
  \item Лидер выполняет атомарное переименование временной директории: \texttt{<model>\_tmp $\to$ <model>}. В POSIX-совместимых файловых системах операция \texttt{rename(2)} является атомарной в пределах одной файловой системы: наблюдатели видят либо отсутствие директории, либо полностью готовую директорию, но никогда не видят промежуточного состояния.
  \item Лидер создаёт маркер-файл \texttt{.ready} в директории модели и освобождает файловую блокировку.
\end{enumerate}

Поды, не получившие эксклюзивную блокировку на шаге~2, устанавливают разделяемую блокировку (\texttt{LOCK\_SH}) на тот же файл \texttt{.lock}. Вызов \texttt{flock(LOCK\_SH)} является блокирующим и возвращает управление только после того, как лидер освободит эксклюзивную блокировку. После возврата из блокирующего вызова под повторно проверяет маркер \texttt{.ready}. Если маркер присутствует, pod переходит к запуску runtime-контейнера. Если маркер отсутствует (например, лидер завершился аварийно до создания маркера), под повторяет попытку получить эксклюзивную блокировку и при успехе становится новым лидером.

Описанный алгоритм в псевдокоде:

\begin{lstlisting}[language=Python, caption={Псевдокод координированной загрузки модели}]
function EnsureModelReady(model_dir, lock_path, ready_marker):
    if Exists(ready_marker):
        return OK

    lock_fd = Open(lock_path, CREATE | RDWR)

    if TryLock(lock_fd, LOCK_EX | LOCK_NB) == SUCCESS:
        # Повторная проверка после получения блокировки
        if Exists(ready_marker):
            Unlock(lock_fd)
            return OK

        tmp_dir = model_dir + "_tmp"
        err = DownloadModel(source_url, tmp_dir)
        if err != nil:
            Unlock(lock_fd)
            RemoveAll(tmp_dir)
            return ERROR

        Rename(tmp_dir, model_dir)   # атомарно
        CreateFile(ready_marker)
        Unlock(lock_fd)
        return OK
    else:
        # Ожидаем завершения лидера
        Lock(lock_fd, LOCK_SH)       # блокирующий вызов
        Unlock(lock_fd)
        # Рекурсивная проверка маркера
        return EnsureModelReady(model_dir, lock_path, ready_marker)
\end{lstlisting}

\subsection{Файловые блокировки flock и их особенности поверх NFS}

\subsubsection*{Advisory locking: разделяемые и эксклюзивные блокировки}

Системный вызов \texttt{flock(2)} реализует рекомендательные (advisory) блокировки на уровне открытого файлового дескриптора \cite{flock_man}. Рекомендательный характер означает, что ядро операционной системы не препятствует доступу к файлу процессам, которые не запрашивают блокировку. Консистентность обеспечивается только при условии, что все участники взаимодействия явно запрашивают и соблюдают соглашение о блокировках.

Поддерживаются два вида блокировок:
\begin{itemize}[leftmargin=*]
  \item \textbf{Разделяемая (shared, \texttt{LOCK\_SH})} --- несколько процессов могут одновременно удерживать разделяемую блокировку на одном файле. Используется для операций чтения, не порождающих конфликтов при параллельном доступе.
  \item \textbf{Эксклюзивная (exclusive, \texttt{LOCK\_EX})} --- только один процесс может удерживать эксклюзивную блокировку. Пока она активна, ни один другой процесс не может получить ни разделяемую, ни эксклюзивную блокировку на тот же файл.
\end{itemize}

Блокировки \texttt{flock} привязаны к описанию открытого файла (open file description), а не к иноду. Это означает, что дублирование файлового дескриптора через \texttt{fork(2)} или \texttt{dup(2)} приводит к тому, что оба дескриптора ссылаются на одно описание и, следовательно, на одну и ту же блокировку. Блокировка автоматически снимается при закрытии последнего файлового дескриптора, ссылающегося на данное описание открытого файла, или при завершении процесса. Это свойство является встроенной защитой от «зависших» блокировок: аварийное завершение процесса-лидера автоматически освобождает блокировку.

\subsubsection*{NFSv3: блокировки через Network Lock Manager}

Поведение \texttt{flock} поверх NFS существенно зависит от версии протокола. В NFSv3 блокировки файлов не являются частью основного протокола. Они реализуются через вспомогательный сервис Network Lock Manager (NLM), который взаимодействует с демонами \texttt{lockd} и \texttt{statd} на стороне клиента и сервера.

Данная архитектура имеет ряд ограничений. Демоны \texttt{lockd} и \texttt{statd} являются дополнительными точками отказа: при их недоступности или некорректной конфигурации операции блокировки возвращают ошибку \texttt{ENOLCK}. При сетевом разделе между клиентом и NFS-сервером блокировка может зависнуть. При перезагрузке NFS-сервера состояние блокировок восстанавливается не всегда корректно, что может привести к ситуации, когда двое клиентов считают себя держателями эксклюзивной блокировки. В ранних версиях ядра Linux (до 2.6.11) \texttt{flock} поверх NFS ограничивался локальной машиной и не распространял блокировку на другие клиенты.

\subsubsection*{NFSv4: встроенные блокировки и lease-механизм}

В NFSv4 файловые блокировки интегрированы непосредственно в основной протокол, что устраняет зависимость от внешних демонов NLM. Сервер NFSv4 отслеживает состояние всех активных блокировок клиентов через механизм аренды (lease): каждый клиент обязан периодически подтверждать свою активность; при истечении аренды сервер освобождает все блокировки данного клиента. Это предотвращает зависшие блокировки даже при аварийном отключении клиента.

Тем не менее NFSv4 не лишён ограничений. В Linux, начиная с версии 2.6.12, клиент NFSv4 эмулирует \texttt{flock} через байт-диапазонные блокировки \texttt{fcntl(F\_SETLK)}, что означает взаимодействие между двумя типами блокировок. Для эксклюзивной блокировки через \texttt{flock} файл должен быть открыт с флагом записи, даже если запись не предполагается. При сетевом разделе (network partition) сервер продолжает удерживать информацию о блокировках до истечения аренды, а не немедленно освобождает их, что создаёт паузу в доступности ресурса.

\subsubsection*{Практическая модель надёжности: flock + rename + маркер}

Анализ ограничений \texttt{flock} поверх NFS показывает, что в производственной среде нельзя полагаться исключительно на корректность блокировок при сетевых сбоях. Именно поэтому описанный алгоритм координированной загрузки использует комбинацию трёх механизмов, образующих эшелонированную защиту:

\begin{enumerate}[leftmargin=*]
  \item \textbf{Файловая блокировка (\texttt{flock})} --- координирует конкурентный доступ в штатном режиме работы, когда сеть стабильна.
  \item \textbf{Атомарное переименование (\texttt{rename})} --- гарантирует, что читатели никогда не видят частично записанных данных. Операция \texttt{rename(2)} атомарна в пределах одной файловой системы согласно POSIX-стандарту.
  \item \textbf{Маркер готовности (\texttt{.ready})} --- является финальным критерием готовности модели. Даже если блокировка была выдана некорректно двум клиентам одновременно, только тот под завершит загрузку, который успеет выполнить \texttt{rename} и создать маркер первым. Второй под, завершив загрузку во временную директорию, обнаружит, что целевая директория уже существует, и примет её как готовую.
\end{enumerate}

Такая комбинация значительно повышает устойчивость алгоритма к нештатным ситуациям, не требуя строгих гарантий от файловой системы.

\subsection{Kubernetes Lease и Leader Election}

\subsubsection*{Ресурс coordination.k8s.io/v1/Lease}

Файловые блокировки поверх NFS предоставляют достаточную надёжность в большинстве сценариев, однако остаются уязвимыми к ситуациям, связанным с некорректной конфигурацией NFS-клиента или нестандартным поведением NFS-сервера. Для случаев, когда требуется более высокий уровень гарантий, Kubernetes предоставляет встроенный механизм распределённой координации через ресурс \texttt{coordination.k8s.io/v1/Lease} \cite{k8s_lease_api}.

Ресурс \texttt{Lease} представляет собой объект Kubernetes API, хранимый в etcd и обладающий строгими гарантиями согласованности за счёт оптимистичной блокировки через поле \texttt{resourceVersion}. Объект содержит следующие ключевые поля спецификации:
\begin{itemize}[leftmargin=*]
  \item \texttt{holderIdentity} --- строковый идентификатор текущего держателя аренды (как правило, имя пода или уникальный идентификатор процесса). Пустое значение означает, что аренда никем не удерживается.
  \item \texttt{leaseDurationSeconds} --- максимальная длительность аренды в секундах без обновления. Если текущий держатель не обновил аренду в течение этого интервала, другие кандидаты вправе захватить её.
  \item \texttt{renewTime} --- временная метка последнего обновления аренды текущим держателем. Лидер обязан периодически обновлять это поле, чтобы подтвердить свою активность.
  \item \texttt{acquireTime} --- временная метка первоначального захвата аренды текущим держателем.
  \item \texttt{leaderTransitions} --- счётчик количества смен лидера, полезный для мониторинга стабильности координации.
\end{itemize}

Kubernetes API-сервер гарантирует, что при одновременной попытке обновить объект \texttt{Lease} несколькими клиентами только один запрос будет успешным: операция обновления проверяет \texttt{resourceVersion}, и если объект был изменён другим клиентом с момента последнего чтения, запрос отклоняется с ошибкой конфликта. Это свойство обеспечивает выбор ровно одного лидера даже при одновременном участии произвольного числа кандидатов.

\subsubsection*{Алгоритм client-go leaderelection}

Стандартная библиотека \texttt{k8s.io/client-go/tools/leaderelection} реализует алгоритм выбора лидера поверх ресурса \texttt{Lease}. Алгоритм работает следующим образом:

\begin{enumerate}[leftmargin=*]
  \item При запуске каждый кандидат пытается создать объект \texttt{Lease} с собственным \texttt{holderIdentity} или обновить существующий объект, если аренда истекла. Операция выполняется через Kubernetes API.
  \item Кандидат, чей запрос на обновление был принят первым (с корректным \texttt{resourceVersion}), становится лидером. Остальные кандидаты получают ответ о конфликте версии и переходят в режим ожидания.
  \item Лидер периодически (с интервалом \texttt{RetryPeriod}) обновляет поле \texttt{renewTime} объекта \texttt{Lease}, не давая сроку аренды истечь. Функция \texttt{OnStartedLeading(ctx)} вызывается при получении статуса лидера и выполняет целевую операцию --- загрузку модели.
  \item Если лидер завершается аварийно или теряет связь с API-сервером, обновление \texttt{renewTime} прекращается. По истечении \texttt{leaseDurationSeconds} аренда помечается как просроченная.
  \item Один из ожидающих кандидатов обнаруживает просроченную аренду, захватывает \texttt{Lease} и становится новым лидером. Функция \texttt{OnStoppedLeading} вызывается на бывшем лидере при обнаружении потери статуса.
\end{enumerate}

Структура \texttt{LeaderElectionConfig} задаёт ключевые параметры: \texttt{LeaseDuration} (типовое значение --- 15 секунд), \texttt{RenewDeadline} (типовое значение --- 10 секунд), \texttt{RetryPeriod} (типовое значение --- 2 секунды), а также набор callback-функций (\texttt{OnStartedLeading}, \texttt{OnStoppedLeading}, \texttt{OnNewLeader}).

\subsubsection*{Координация per-model через пространство имён Lease}

В контексте загрузки моделей ресурс \texttt{Lease} именуется по схеме \texttt{model-<name>-<revision>} в пространстве имён соответствующего сервиса инференса. Это позволяет независимо координировать загрузку разных моделей: одновременно могут присутствовать несколько объектов \texttt{Lease}, каждый из которых управляет загрузкой конкретной версии конкретной модели.

Лидер, захвативший \texttt{Lease} для данной пары (модель, ревизия), выполняет загрузку в NFS, создаёт маркер \texttt{.ready} и освобождает \texttt{Lease}. Остальные поды, работающие с той же моделью, отслеживают либо состояние \texttt{Lease} через механизм Watch, либо появление маркера \texttt{.ready} в NFS-директории, либо используют оба сигнала в conjunction. Данный подход исключает условие гонки на уровне Kubernetes без каких-либо требований к семантике NFS-блокировок.

Преимущество Kubernetes Lease перед \texttt{flock} состоит в том, что гарантии согласованности обеспечиваются etcd --- распределённым хранилищем с алгоритмом Raft, а не файловой системой. Таким образом, Leader Election через Kubernetes Lease является предпочтительным механизмом координации для случаев, когда init-контейнер имеет доступ к Kubernetes API. В случаях, когда доступ к API недоступен или нежелателен (например, в минималистичных окружениях), \texttt{flock} с атомарным переименованием остаётся резервным вариантом.

\subsection{Ограничения NFS-этапа}

\subsubsection*{Централизованный сервер как единая точка отказа}

Несмотря на существенное улучшение по сравнению со схемой \texttt{emptyDir}, NFS-подход обнаружил ряд системных ограничений, которые были выявлены в процессе промышленной эксплуатации и ставшие основанием для разработки следующего этапа архитектуры.

\textbf{NFS-сервер как SPOF.} Весь ввод-вывод чтения весов моделей проходит через один NFS-сервер. При его недоступности --- перезагрузке, сетевой проблеме, исчерпании дискового пространства или переполнении очереди запросов --- вся система инференса лишается доступа к моделям. Kubernetes не содержит встроенных механизмов высокой доступности для внешнего NFS-сервера: сбой сервера воспринимается как продолжительный таймаут сетевой операции, что приводит к зависанию подов на этапе монтирования или чтения.

\textbf{Bottleneck пропускной способности.} При одновременном масштабировании большого числа реплик все они читают файлы модели через сетевую карту NFS-сервера. Пропускная способность его сетевого интерфейса физически ограничена и не масштабируется горизонтально. При burst-масштабировании с 0 до $N$ реплик суммарный входящий трафик на NFS-сервер достигает $N \times \text{скорость чтения}$, что при типичных размерах моделей (7--70~ГБ) быстро приводит к насыщению канала.

\textbf{Отсутствие локальности данных.} Даже если узел уже обслуживал данную модель ранее, при создании нового пода runtime-контейнер читает данные из NFS по сети заново. Локальный диск узла не используется для кэширования: информация, однажды прочитанная с NFS, не сохраняется на узле и при следующем запросе снова загружается по сети. Это приводит к тому, что задержка cold start определяется не скоростью локального чтения (которая может достигать 3--5~ГБ/с для NVMe), а пропускной способностью сети между нодой и NFS-сервером.

\textbf{Сложность построения отказоустойчивого NFS.} Создание высокодоступного NFS-сервера требует значительных дополнительных усилий: кластеризации с использованием Pacemaker/Corosync, синхронной репликации данных через DRBD или аналогичные средства, настройки shared-storage для кластерного NFS. Такое решение существенно увеличивает операционную сложность и стоимость инфраструктуры, при этом не устраняя принципиальных проблем с масштабируемостью полосы пропускания.

\textbf{Нагрузка на сеть при burst-масштабировании.} В отличие от схемы \texttt{emptyDir}, где нагрузка распределялась между нодами (каждый узел загружал из внешнего источника независимо), в NFS-схеме весь трафик концентрируется на одном сервере. При резком росте числа реплик весь кластер конкурирует за пропускную способность одного сетевого пути.

\subsubsection*{CAP-теорема и выбор NFS}

CAP-теорема \cite{brewer2000cap} утверждает, что распределённая система не может одновременно обеспечивать согласованность (Consistency), доступность (Availability) и устойчивость к разделению сети (Partition tolerance). При разрыве сети (Partition) между нодами Kubernetes и NFS-сервером система вынуждена выбирать между двумя оставшимися свойствами.

Классический NFS с режимом монтирования \texttt{hard} выбирает согласованность, жертвуя доступностью: при недоступности сервера все операции ввода-вывода блокируются и ожидают восстановления соединения. Runtime-контейнер, ожидающий чтения весов из NFS, зависает до тех пор, пока сетевой сбой не будет устранён. Режим монтирования \texttt{soft} позволяет операциям завершаться с ошибкой по таймауту, однако приводит к потенциальной несогласованности данных при восстановлении.

В итоге NFS-этап предлагает приемлемые характеристики для небольших кластеров с умеренной нагрузкой, но системно неприменим в условиях масштабирования, требующего высокой параллельности и устойчивости к сбоям сети.

\subsection{Выводы по разделу}

NFS-этап значительно сократил количество обращений к внешнему репозиторию: вместо $N$ независимых загрузок на $N$ реплик однократная загрузка в общее хранилище обеспечивает повторное использование данных всеми подами. Алгоритм координации через файловые блокировки \texttt{flock} и маркер готовности обеспечивает консистентность при параллельном старте. Kubernetes Lease API предоставляет более надёжную альтернативу для координации в случаях, когда гарантии NFS-блокировок недостаточны.

Вместе с тем NFS как централизованное сетевое хранилище переносит проблему с внешнего репозитория на один внутренний сервер, создавая единую точку отказа и ограниченный по пропускной способности ресурс. При масштабировании до десятков параллельных реплик, каждая из которых читает модель размером 7--70~ГБ, NFS-сервер становится узким местом архитектуры. Отсутствие локального кэширования на дисках нод означает, что одни и те же данные многократно передаются по сети при каждом событии масштабирования.

Для устранения этих ограничений необходим переход к архитектуре, в которой данные хранятся на локальных дисках worker-нод и доставляются напрямую между агентами, минуя централизованный сервер. Такая архитектура, реализованная через распределённый кэш на уровне нод с P2P-взаимодействием и FUSE-клиентом, обеспечивает локальность данных, горизонтальную масштабируемость пропускной способности и устойчивость к отказу отдельных компонентов. Именно она составляет содержание следующего раздела.
